\chapter{L2 -- Sistemi a stato vettore, rappresentazioni, moto, equilibrio, raggiungibilità, osservabilità e sistemi lineari}

\section{Sistemi a stato vettore}

In questa lezione si passa dalla descrizione astratta del sistema dinamico alla formulazione \emph{a stato vettore}.  
L'idea è che, in moltissimi problemi di interesse ingegneristico, i valori istantanei di ingresso, uscita e stato possano essere rappresentati come vettori.

Si assumono quindi:
\[
U,\quad Y,\quad X
\]
spazi vettoriali, con
\[
\dim U = m,\qquad \dim Y = \ell,\qquad \dim X = n.
\]

Qui:
\begin{itemize}
    \item $m$ è il numero di ingressi;
    \item $\ell$ è il numero di uscite;
    \item $n$ è il numero di variabili di stato.
\end{itemize}

\paragraph{Caso SISO e MIMO.}
Se
\[
m=\ell=1,
\]
il sistema è SISO (\emph{Single Input Single Output}); altrimenti, in generale, è MIMO.

\subsection{Spazi dei valori istantanei e spazi dei segnali}

Anche qui è fondamentale distinguere:
\[
U \ \text{da}\ \mathcal{U},
\qquad
Y \ \text{da}\ \mathcal{Y}.
\]

Infatti:
\begin{itemize}
    \item $U$ e $Y$ sono spazi vettoriali di valori \emph{istantanei};
    \item $\mathcal{U}$ e $\mathcal{Y}$ sono insiemi di \emph{funzioni nel tempo}.
\end{itemize}

Se $u_1(\cdot),u_2(\cdot)\in\mathcal{U}$ e $\alpha\in\mathbb{R}$, le operazioni naturali sono definite punto per punto:
\[
(u_1+u_2)(t)=u_1(t)+u_2(t),
\qquad
(\alpha u_1)(t)=\alpha\,u_1(t).
\]

Quando $\mathcal{U}$ è chiuso rispetto a tali operazioni, esso è uno spazio vettoriale. Lo stesso vale per $\mathcal{Y}$.

\paragraph{Osservazione importante.}
In generale:
\begin{itemize}
    \item $U$, $X$, $Y$ sono spazi vettoriali di dimensione finita;
    \item $\mathcal{U}$ e $\mathcal{Y}$ sono invece, tipicamente, spazi di dimensione infinita.
\end{itemize}

Questo diventerà cruciale quando si parlerà di rappresentazioni matriciali: una mappa lineare tra spazi finito-dimensionali si rappresenta con una matrice; una mappa lineare definita su uno spazio di funzioni, in generale, no.

\section{Rappresentazione esplicita del sistema}

Nel linguaggio introdotto nella lezione precedente, un sistema dinamico può essere descritto in forma esplicita come
\begin{align}
x(t) &= \phi\bigl(t,t_0,x(t_0),u_{[t_0,t]}\bigr), \label{eq:L2_phi}\\
y(t) &= g\bigl(t,t_0,x(t_0),u_{[t_0,t]}\bigr). \label{eq:L2_gglob}
\end{align}

Questa scrittura è detta \emph{esplicita} perché fornisce direttamente lo stato e l'uscita al tempo $t$.

\paragraph{Interpretazione.}
Se conosco:
\begin{itemize}
    \item lo stato iniziale $x(t_0)$,
    \item l'ingresso applicato nell'intervallo $[t_0,t]$,
    \item la funzione di transizione $\phi$,
\end{itemize}
allora conosco esattamente lo stato al tempo $t$.

\section{Rappresentazione implicita}

Spesso il sistema non viene descritto direttamente tramite $\phi$, ma tramite una legge locale di evoluzione.

\subsection{Caso a tempo continuo}

Nel tempo continuo, sotto opportune ipotesi di regolarità, la descrizione diventa:
\begin{equation}
\dot x(t)=f\bigl(x(t),u(t),t\bigr),
\label{eq:L2_ct_state}
\end{equation}
con condizione iniziale
\[
x(t_0)=x_0.
\]

L'uscita è descritta localmente da
\begin{equation}
y(t)=g\bigl(x(t),u(t),t\bigr).
\label{eq:L2_ct_output}
\end{equation}

Questa è la \emph{rappresentazione implicita}: non mi dice direttamente quale sia lo stato a un certo tempo futuro, ma mi dice \emph{come varia} lo stato.

\subsection{Da \texorpdfstring{$\phi$}{phi} a \texorpdfstring{$f$}{f}: idea del passaggio}

Dalla proprietà di composizione si ha, per $\Delta t>0$ piccolo,
\[
x(t+\Delta t)=\phi\bigl(t+\Delta t,t,x(t),u_{[t,t+\Delta t]}\bigr).
\]

Allora
\[
\frac{x(t+\Delta t)-x(t)}{\Delta t}
=
\frac{\phi\bigl(t+\Delta t,t,x(t),u_{[t,t+\Delta t]}\bigr)-x(t)}{\Delta t}.
\]

Se il limite esiste per $\Delta t\to 0$, si ottiene
\[
\dot x(t)=f\bigl(x(t),u(t),t\bigr).
\]

\paragraph{Lettura concettuale.}
\begin{itemize}
    \item La forma esplicita descrive il risultato finale dell'evoluzione.
    \item La forma implicita descrive la legge istantanea che genera quell'evoluzione.
\end{itemize}

\subsection{Caso a tempo discreto}

Nel tempo discreto non si usa la derivata, ma una relazione ricorsiva:
\begin{equation}
x(k+1)=f\bigl(x(k),u(k),k\bigr),
\label{eq:L2_dt_state}
\end{equation}
e
\begin{equation}
y(k)=g\bigl(x(k),u(k),k\bigr).
\label{eq:L2_dt_output}
\end{equation}

Qui la rappresentazione implicita coincide di fatto con un aggiornamento passo-passo dello stato.

\paragraph{Osservazione.}
Nel continuo passare da $f$ a $\phi$ significa risolvere un sistema di equazioni differenziali.  
Nel discreto passare da $f$ a $\phi$ significa iterare la dinamica.

\section{Eventi, moti, traiettorie e punti di equilibrio}

\subsection{Evento}

Un \emph{evento} del sistema è una coppia
\[
(x(t),t)\in X\times T.
\]

Si tratta quindi di uno stato osservato in un preciso istante.

\begin{figure}[H]
\centering
\begin{tikzpicture}[scale=1.0]
    \draw[->] (0,0) -- (7.5,0) node[right] {$t$};
    \draw[->] (0,0) -- (0,4.0) node[above] {$x(t)$};

    \draw[thick]
        plot[smooth] coordinates {(0.6,1.0) (1.3,1.8) (2.2,1.5) (3.2,2.4) (4.1,2.0) (5.3,3.0) (6.5,3.4)};

    \draw[dashed] (3.2,0) -- (3.2,2.4);
    \draw[dashed] (0,2.4) -- (3.2,2.4);
    \filldraw (3.2,2.4) circle (2pt);

    \node at (3.55,2.75) {evento $(x(t),t)$};
\end{tikzpicture}
\caption{Rappresentazione grafica di un evento del sistema.}
\end{figure}

\subsection{Moto del sistema}

Fissati:
\begin{itemize}
    \item uno stato iniziale $x(t_0)$,
    \item un ingresso $u(\cdot)$,
    \item un istante iniziale $t_0$,
\end{itemize}
il \emph{moto} del sistema è l'insieme
\[
\Bigl\{(x(t),t)\in X\times T:\ x(t)=\phi(t,t_0,x(t_0),u),\ t\ge t_0\Bigr\}.
\]

Il moto vive quindi nello spazio stato--tempo.

\subsection{Traiettoria}

La \emph{traiettoria} è la proiezione del moto sul solo spazio degli stati:
\[
\Bigl\{x(t)\in X:\ x(t)=\phi(t,t_0,x(t_0),u),\ t\ge t_0\Bigr\}.
\]

\paragraph{Differenza tra moto e traiettoria.}
\begin{itemize}
    \item Il moto tiene conto anche del tempo.
    \item La traiettoria descrive solo il percorso dello stato nello spazio $X$.
\end{itemize}

\subsection{Punto di equilibrio}

Uno stato $\bar x\in X$ è un \emph{punto di equilibrio} per un ingresso $\bar u(\cdot)$ se
\[
\bar x=\phi(t,t_0,\bar x,\bar u)\qquad \forall t\ge t_0.
\]

Questo significa che, se il sistema parte da $\bar x$ ed è soggetto all'ingresso $\bar u$, allora resta per sempre nello stesso stato.

\paragraph{Tempo continuo.}
Se il sistema è descritto da
\[
\dot x(t)=f(x(t),u(t),t),
\]
allora la condizione di equilibrio è
\[
0=f(\bar x,\bar u(t),t)\qquad \forall t\ge t_0.
\]

Nel caso tempo-invariante con ingresso costante $\bar u$, si riduce a
\[
0=f(\bar x,\bar u).
\]

\paragraph{Tempo discreto.}
Se
\[
x(k+1)=f(x(k),u(k),k),
\]
allora un equilibrio soddisfa
\[
\bar x=f(\bar x,\bar u(k),k)\qquad \forall k\ge k_0.
\]

Nel caso tempo-invariante con ingresso costante:
\[
\bar x=f(\bar x,\bar u).
\]

\section{Ciclo limite}

Un moto è periodico se esiste $T_p>0$ tale che
\[
x(t+T_p)=x(t)\qquad \forall t.
\]

Un \emph{ciclo limite} è una traiettoria periodica isolata nello spazio di stato.

\paragraph{Interpretazione.}
Il sistema non converge a un punto di equilibrio, ma evolve ripetutamente lungo una stessa orbita chiusa.

\begin{figure}[H]
\centering
\begin{tikzpicture}[scale=1.0]
    \draw[->] (-3.2,0) -- (3.2,0) node[right] {$x_1$};
    \draw[->] (0,-2.7) -- (0,2.7) node[above] {$x_2$};

    \draw[thick]
        plot[smooth cycle, tension=0.9]
        coordinates {(1.8,0.2) (1.0,1.6) (-0.8,1.7) (-1.8,0.2) (-1.2,-1.5) (0.8,-1.7)};
\end{tikzpicture}
\caption{Schema qualitativo di un ciclo limite nel piano di stato.}
\end{figure}

\paragraph{Nota concettuale.}
Il ciclo limite è tipico dei sistemi non lineari. Nei sistemi lineari possono comparire moti periodici, ma non nel senso strutturale forte tipico dei cicli limite isolati.

\section{Raggiungibilità e controllabilità}

L'ingresso modifica il moto del sistema. Da qui nascono due nozioni fondamentali.

Siano dati due eventi:
\[
(x_1,t_1),\qquad (x_2,t_2),\qquad t_2>t_1.
\]

Si dice che $(x_2,t_2)$ è \emph{raggiungibile} da $(x_1,t_1)$ se esiste un ingresso $u(\cdot)\in\mathcal{U}$ tale che
\[
x_2=\phi(t_2,t_1,x_1,u).
\]

In modo duale, si dice che $(x_1,t_1)$ è \emph{controllabile} da $(x_2,t_2)$ se esiste un ingresso che porta il sistema da $x_1$ al tempo $t_1$ a $x_2$ al tempo $t_2$.

\paragraph{Idea intuitiva.}
\begin{itemize}
    \item Raggiungibilità: posso arrivare nello stato desiderato?
    \item Controllabilità: posso forzare il sistema a seguire l'evoluzione voluta?
\end{itemize}

\section{Osservabilità e ricostruibilità}

La domanda di fondo è:

\begin{quote}
Osservando ingresso e uscita, riesco a capire da quale stato iniziale è partito il sistema?
\end{quote}

\subsection{Stati indistinguibili}

Due stati iniziali $x_0$ e $\bar x_0$ si dicono \emph{indistinguibili} nell'intervallo $[t_0,t_1]$ se, applicando lo stesso ingresso, producono la stessa uscita su tutto l'intervallo:
\[
g(\tau,t_0,x_0,u)=g(\tau,t_0,\bar x_0,u)
\qquad
\forall \tau\in [t_0,t_1],\ \forall u\in\mathcal U.
\]

Se esistono due stati distinti indistinguibili, il sistema non è osservabile.

\paragraph{Interpretazione fisica.}
Il sensore non riesce a distinguere quei due stati: l'uscita ``copre'' troppo poco dello stato.

\section{Esempio discreto di mancata raggiungibilità e mancata osservabilità}

Consideriamo il sistema a tempo discreto
\[
\begin{cases}
x(k+1)=Ax(k)+Bu(k),\\
y(k)=Cx(k),
\end{cases}
\]
con
\[
A=
\begin{bmatrix}
0 & 1\\
0 & 0
\end{bmatrix},
\qquad
B=
\begin{bmatrix}
1\\
0
\end{bmatrix},
\qquad
C=
\begin{bmatrix}
0 & 1
\end{bmatrix}.
\]

\subsection{Verifica della raggiungibilità}

La matrice di raggiungibilità è
\[
\mathcal{R}=\begin{bmatrix}B & AB\end{bmatrix}.
\]

Calcoliamo:
\[
AB=
\begin{bmatrix}
0 & 1\\
0 & 0
\end{bmatrix}
\begin{bmatrix}
1\\
0
\end{bmatrix}
=
\begin{bmatrix}
0\\
0
\end{bmatrix}.
\]

Quindi
\[
\mathcal{R}=
\begin{bmatrix}
1 & 0\\
0 & 0
\end{bmatrix},
\qquad
\operatorname{rank}(\mathcal{R})=1<2.
\]

Il sistema non è raggiungibile.

\paragraph{Lettura diretta.}
Se parto da
\[
x(0)=\begin{bmatrix}0\\0\end{bmatrix},
\]
allora
\[
x(1)=Ax(0)+Bu(0)=\begin{bmatrix}u(0)\\0\end{bmatrix},
\]
e poi
\[
x(2)=Ax(1)+Bu(1)=\begin{bmatrix}u(1)\\0\end{bmatrix}.
\]

La seconda componente resta sempre nulla.  
Quindi non posso raggiungere uno stato del tipo
\[
\begin{bmatrix}
1\\5
\end{bmatrix}.
\]

\subsection{Verifica dell'osservabilità}

La matrice di osservabilità è
\[
\mathcal{O}=
\begin{bmatrix}
C\\
CA
\end{bmatrix}.
\]

Calcoliamo:
\[
CA=
\begin{bmatrix}
0 & 1
\end{bmatrix}
\begin{bmatrix}
0 & 1\\
0 & 0
\end{bmatrix}
=
\begin{bmatrix}
0 & 0
\end{bmatrix}.
\]

Quindi
\[
\mathcal{O}=
\begin{bmatrix}
0 & 1\\
0 & 0
\end{bmatrix},
\qquad
\operatorname{rank}(\mathcal{O})=1<2.
\]

Il sistema non è osservabile.

\paragraph{Interpretazione.}
L'uscita vale
\[
y(k)=Cx(k)=
\begin{bmatrix}
0 & 1
\end{bmatrix}
\begin{bmatrix}
x_1(k)\\
x_2(k)
\end{bmatrix}
=x_2(k).
\]

Il sensore legge solo la seconda componente dello stato.  
La prima componente può cambiare senza essere direttamente rilevata: per questo parte dello stato rimane nascosta all'osservazione.

\section{Sistemi lineari}

Un sistema si dice lineare se la funzione di transizione $\phi$ e la trasformazione d'uscita $g$ sono lineari rispetto a stato iniziale e ingresso.

Formalmente, per ogni $x_1,x_2\in X$, $u_1,u_2\in\mathcal{U}$ e $\alpha_1,\alpha_2\in\mathbb{R}$,
\begin{multline}
\phi\bigl(t,t_0,\alpha_1x_1+\alpha_2x_2,\alpha_1u_1+\alpha_2u_2\bigr)
\\
=
\alpha_1\phi(t,t_0,x_1,u_1)+\alpha_2\phi(t,t_0,x_2,u_2),
\label{eq:L2_linearity_phi}
\end{multline}
e analogamente
\begin{multline}
g\bigl(t,t_0,\alpha_1x_1+\alpha_2x_2,\alpha_1u_1+\alpha_2u_2\bigr)
\\
=
\alpha_1g(t,t_0,x_1,u_1)+\alpha_2g(t,t_0,x_2,u_2).
\label{eq:L2_linearity_g}
\end{multline}

\subsection{Proprietà di decomposizione}

La linearità permette di separare l'effetto dello stato iniziale dall'effetto dell'ingresso.

Ponendo in \eqref{eq:L2_linearity_phi}
\[
x_2=0,\qquad u_1=0,\qquad \alpha_1=\alpha_2=1,
\]
si ottiene
\[
\phi(t,t_0,x,u)=\phi(t,t_0,x,0)+\phi(t,t_0,0,u).
\]

Quindi l'evoluzione dello stato si decompone in:
\[
\phi(t,t_0,x,u)=\phi_\ell(t,t_0,x)+\phi_f(t,t_0,u),
\]
dove
\[
\phi_\ell(t,t_0,x):=\phi(t,t_0,x,0)
\]
è la \emph{risposta libera}, e
\[
\phi_f(t,t_0,u):=\phi(t,t_0,0,u)
\]
è la \emph{risposta forzata}.

Allo stesso modo:
\[
g(t,t_0,x,u)=g_\ell(t,t_0,x)+g_f(t,t_0,u).
\]

\paragraph{Interpretazione.}
\begin{itemize}
    \item La risposta libera è quella dovuta unicamente allo stato iniziale.
    \item La risposta forzata è quella dovuta unicamente all'ingresso.
\end{itemize}

\begin{figure}[H]
\centering
\begin{tikzpicture}[>=Latex, scale=1.0]
    \node[draw, rounded corners, minimum width=2.7cm, minimum height=1.1cm] (sum) at (0,0) {$\phi$};
    \draw[->, thick] (-3,0.7) -- (sum.west) node[midway, above] {$x(t_0)$};
    \draw[->, thick] (-3,-0.7) -- (sum.west) node[midway, below] {$u(\cdot)$};
    \draw[->, thick] (sum.east) -- (3.5,0) node[midway, above] {$x(t)$};

    \node at (0,-2.0) {$x(t)=\underbrace{\phi_\ell(t,t_0,x(t_0))}_{\text{risposta libera}}
    +\underbrace{\phi_f(t,t_0,u)}_{\text{risposta forzata}}$};
\end{tikzpicture}
\caption{Idea di decomposizione della risposta nei sistemi lineari.}
\end{figure}

\subsection{Caso lineare tempo-invariante continuo}

Per un sistema LTI continuo
\[
\dot x(t)=Ax(t)+Bu(t),
\qquad
y(t)=Cx(t)+Du(t),
\]
la soluzione dello stato è
\[
x(t)=e^{A(t-t_0)}x(t_0)+\int_{t_0}^{t} e^{A(t-\tau)}Bu(\tau)\,d\tau.
\]

Qui:
\begin{itemize}
    \item
    \[
    e^{A(t-t_0)}x(t_0)
    \]
    è la risposta libera;
    \item
    \[
    \int_{t_0}^{t} e^{A(t-\tau)}Bu(\tau)\,d\tau
    \]
    è la risposta forzata.
\end{itemize}

\section{Nota finale sulle mappe lineari}

La decomposizione mette in evidenza una differenza importante.

\subsection{Parte libera}

La mappa
\[
\phi_\ell: X\to X
\]
è una mappa lineare tra spazi finito-dimensionali; quindi si può rappresentare con una matrice.

\subsection{Parte forzata}

La mappa
\[
\phi_f:\mathcal{U}\to X
\]
è una mappa lineare da uno spazio di funzioni, in generale infinito-dimensionale, verso lo spazio di stato.

Per questo motivo, in generale, $\phi_f$ \emph{non} si rappresenta con una matrice finita.

\paragraph{Conseguenza.}
Nei sistemi lineari:
\begin{itemize}
    \item la dinamica interna si presta naturalmente a una descrizione matriciale;
    \item l'effetto dell'ingresso richiede invece operatori più generali, oppure formule integrali/sommatorie.
\end{itemize}

\section{Sintesi della lezione}

In questa seconda lezione sono emersi i seguenti punti chiave:
\begin{enumerate}
    \item nei sistemi a stato vettore, ingresso, uscita e stato vivono in spazi vettoriali;
    \item la rappresentazione esplicita usa $\phi$, quella implicita usa $f$;
    \item si definiscono evento, moto, traiettoria, equilibrio e ciclo limite;
    \item si introducono raggiungibilità e osservabilità;
    \item nei sistemi lineari vale la decomposizione in risposta libera e risposta forzata.
\end{enumerate}

Questi concetti saranno alla base di tutto il seguito del corso: esponenziale di matrice, stabilità, raggiungibilità, osservabilità e progetto del controllo.